%-------------------GENERAL--------------------
\usepackage{amsmath,latexsym, amssymb, dsfont}
\usepackage{braket}

%-------------------UNITS--------------------
\usepackage{siunitx}
\DeclareSIUnit{\Mpc}{Mpc}
\DeclareSIUnit{\solarMass}{\ensuremath{M_\odot}}
\DeclareSIUnit{\erg}{erg}
\DeclareSIUnit{\year}{yr}
\DeclareSIUnit{\nats}{nats}

%------------------SHORTCUTS-------------------

\newcommand{\mean}{\overset{-}{T}}
\newcommand{\degree}{^{\circ}}

\newcommand{\inv}{^{-1}}

\newcommand{\sinc}{\text{sinc}}
\newcommand{\const}{\text{const}}

\newcommand{\Lp}{\text{L}}
\newcommand{\Hilbert}{\mathcal{H}}

\newcommand{\N}{\mathbb{N}}
\newcommand{\Z}{\mathbb{Z}}
\newcommand{\R}{\mathbb{R}}
\newcommand{\C}{\mathbb{C}}

\newcommand{\0}{\hat{\mathbf{0}}}
\newcommand{\1}{\hat{\mathds{1}}}

% Inner and outer product
\newcommand{\ip}[2]{\bra{#1}\!\!\ket{#2}}
\newcommand{\op}[2]{\ket{#1}\!\!\bra{#2}}

\newcommand{\norm}[1]{\lVert #1 \rVert}

%-------------MATRIX ENVIRONMENTS-------------

% Three matrix types [], (), ||, chosen by options b, p, v. Standard ist b.
\newcommand{\bmat}[1]{\begin{bmatrix} #1 \end{bmatrix}}
\newcommand{\pmat}[1]{\begin{pmatrix} #1 \end{pmatrix}}
\newcommand{\vmat}[1]{\begin{vmatrix} #1 \end{vmatrix}}
%-------------THEOREM ENVIRONMENTS-------------

%Hier bei Bedarf Definitionen der Theorem-Umgebungen.